\section{SF-E846: homogeneous quadratic point modifications}\label{SF-E846:section}

\textbf{Checked mathematical status, 2026-11-22.} The author, root,
and independent reviewer \texttt{/root/g34\_eightfolds} have read the complete
frozen proof; all three mathematical reviews passed, with no repair.
The independent full proof/source/code report was completed by
\texttt{2026-11-22 02:30:14 UTC}. The frozen Markdown original retains its
author-review status. The author and root have also read all 604 lines of the canonical conversion; all 19 frozen display bodies are byte-identical. No strict comparison with all previous local Artin exclusions, and no independent main-gate movement, is claimed.
Author: \texttt{/root/sixfolds}, GPT-6 Astra Pro; inherited reasoning effort is not exposed.
The all-order noncommutative semiregularity implication is not assumed here.
The global argument below instead constructs a class in ordinary Ext.
Its exact scope is homogeneous scalar local extensions; arbitrary filtered
perturbations and arbitrary matrix-valued extension classes are not included.

\subsection{1. The exact local module and its resolution}\label{SF-E846:s1}

Put \(R=\mathbb C[[x_1,x_2,x_3,x_4]]\). Let \(f_1,\ldots,f_4\) be homogeneous quadrics forming a
regular sequence, \(I=(f_1,\ldots,f_4)\), and \(A=R/I\). Let \(U=\mathbb C^4\) with ordered basis
\(e_1,\ldots,e_4\). Let \(\iota_i\) denote contraction by the dual basis element and let
\(e_i\wedge\) denote exterior multiplication. Write

\[
 K_p=R\otimes\Lambda^pU,\quad d=\sum_i f_i\iota_i,
 \qquad
 \beta=\sum_{i<j}b_{ij}\iota_j\iota_i,
 \tag{HK.1}
\]

where the six \(b_{ij}\) are linear forms, not all zero. Scalars commute
throughout. No spanning hypothesis on the \(b_{ij}\) is imposed. The convention makes
\(\beta(e_i\wedge e_j)=b_{ij}\) for \(i<j\). The even contraction \(\beta\) commutes with \(d\).
There is an actual extension \(M_b\) with graded exact sequence

\[
 0\longrightarrow A(-3)\longrightarrow M_b\longrightarrow I
 \longrightarrow0.
 \tag{HK.2}
\]

Here an \(R(-j)\) generator has degree \(j\). A resolution has
\(L_p=K_p(-3)\) direct-sum \(K_{p+1}\) for \(0\le p\le3\), with \(L_4=K_4(-3)\).
In this notation the first \(K_p\) generator has degree \(3+2p\) and the second
\(K_{p+1}\) generator has degree \(2+2p\). Its differential is

\[
 D=\begin{pmatrix}d&\beta\\0&-d\end{pmatrix}.
 \tag{HK.3}
\]

More explicitly the bottom complex is the resolution of \(I\), ending in \(K_1\);
the top complex resolves \(A(-3)\). Their triangular extension is exact in
negative cohomological degrees by the cohomology sequence of the two
complexes. Its degree-zero cohomology is (HK.2). The canonical hull map
\(\phi:M_b\to R\) is the map onto \(I\) followed by inclusion. Its kernel is \(A(-3)\)
and its cokernel is \(A\). The module is rank one and is locally free away
from the closed point. All displayed maps are continuous for the maximal
ideal topology. The notation \(K_p(-3)\) is shorthand for the stated generator
degrees, not for a second independent shift of an already shifted \(K_p\).

\subsection{2. First-order corrections, with every lowest weight exhausted}\label{SF-E846:s2}

Let \(\pi\) be a constant nondegenerate bivector. Use right modules and the
Moyal product

\[
 a\star b=ab+\frac h2\{a,b\}_\pi+
 \frac{h^2}{8}\sum\pi^{ij}\pi^{kl}
       (\partial_i\partial_k a)(\partial_j\partial_l b)
        \pmod {h^3}.
 \tag{HK.4}
\]

The opposite convention is obtained by replacing \(\pi\) by \(-\pi\); existence
and nonexistence statements are unchanged. A matrix acts on a column
of a free right module by left multiplication. Composition of matrices
therefore uses star products in the displayed order.

Give \(h\) degree two. Assume \(M_b\) admits a flat first-order deformation.
The finite free resolution lifts with the same ranks, by successive
kernel lifting and nilpotent Nakayama exactly as in E819 section 2.1.
Write its differential \(D+hH\). A homogeneous endomorphism of cohomological
degree one and internal degree below -3 is zero by the generator shifts.
The degree -3 part can only be a constant lower-left block \(c_p:K_p\to K_p\)
(\(1\le p\le4\)). The equation \(dc-cd=0\) implies \(c_p=C_0\) identity on the exterior
factor. To see this start with \(K_4\), which has exterior rank one, then
compare the coefficients of the four independent quadrics successively
on \(K_4\),\(K_3\),\(K_2\). Contractions of \(K_4\) span \(K_3\), of \(K_3\) span \(K_2\), and of \(K_2\) span \(K_1\).
The diagonal equations also give \(\beta C_0=C_0\beta=0\). Since some \(b_{ij}\) is
nonzero, \(C_0=0\). Thus the entire degree -3 correction vanishes.

The homogeneous part \(H\) of degree -2 has the form

\[
 H=\begin{pmatrix}a&0\\C&e\end{pmatrix},
 \tag{HK.5}
\]

where \(a,e\) have constant coefficients and exterior degree -1, while \(C\)
is a linear scalar acting identically on each exterior degree. Indeed
the same chain argument now compares linear multiples of the \(f_i\).
There are no nonzero linear syzygies among a regular sequence of quadrics:
its Koszul first syzygies start in total degree four. This forces the
lower-left operators to be \(C\) times exterior identity. Upper-right degree
-2 corrections would need a coefficient of degree -1 and hence are zero.

Put \(P_{ij}=\{f_i,f_j\}_\pi\). The diagonal equations are

\[
 \{d,a\}=\iota_{P-Cb},\qquad
 \{d,e\}=-\iota_{P-Cb}.
 \tag{HK.6}
\]

In particular there are unique constant two-forms \(A_l\) on \(U\) such that
\(P-Cb=\sum_l f_l A_l\). Define \(h_A=\sum_l(e_l\wedge)\iota_{A_l}\). Then
\(\{d,h_A\}=\iota_{P-Cb}\). All constant solutions of the first equation in (HK.6)
are \(h_A+\iota_v\), for \(v\) in \(U^*\). The second equation gives \(-h_A+\iota_w\) on its
truncated domain \(K_2\),\(K_3\),\(K_4\). Here is the complete uniqueness argument.
After subtracting \(h_A\), comparison of the coefficients \(f_i\) says the
operator anticommutes with every \(\iota_i\). Its value on \(K_4\) determines a
unique contraction \(\iota_v\). After subtracting that contraction, the
equations on \(K_4\) force the operator to vanish on \(K_3\), those on \(K_3\) force
vanishing on \(K_2\), and those on \(K_2\) force vanishing on \(K_1\). For the bottom
truncation only \(K_4\) and \(K_3\) are needed to determine \(K_3\) and \(K_2\); no equation
on an absent \(K_1\to K_0\) bottom map is used. Thus

\[
 a=h_A+\iota_v,\qquad e=-h_A+\iota_w.
 \tag{HK.7}
\]

For reference the remaining first-order equation is the genuine
off-diagonal identity

\[
 G_\pi+[h_A,\beta]+\iota_{v+w}\beta=0,
 \qquad
 G_\pi=\sum\pi^{ij}(\partial_i d)(\partial_j\beta).
 \tag{HK.8}
\]

It is not replaced by scalar cross minors. No existence of a solution
of (HK.8) for a previously unexcluded two-direction example is claimed.

\subsection{3. A nonzero second-order obstruction of internal degree -4}\label{SF-E846:s3}

First \(C\) is nonzero. If \(C=0\), (HK.6) says every \(P_{ij}\) vanishes in \(A\). The
Pfaffian identity for the nondegenerate \(\pi\) says \(\operatorname{Pf}(P)\) is a nonzero
constant times the Jacobian determinant of \(f\). The latter is nonzero
in \(A\) by the characteristic-zero complete-intersection socle theorem,
contradiction. This is the already checked Eisenbud--Ulrich Theorem 7.1,
p.~210, with full Appendix A.1/A.3/A.5 proofs, pp.~222--225.\cite{SFEisenbudUlrich2019}

Take the degree -2 solution \(H\) furnished by a first-order lift.
The coefficient of \(h^2\) in (\(D+hH\)) star (\(D+hH\)) is a cocycle \(O\) of
cohomological degree two and internal degree -4. The lower-left block
is particularly simple:

\[
 O_{21}=-\iota_{\{f,C\}_\pi}+C\iota_{v+w}.
 \tag{HK.9}
\]

Indeed the second-order term involving \(D\) twice has zero lower-left
block. The two first-order star terms with \(D\),\(H\) give
\(-\sum\pi^{ij}(\partial_i d)(\partial_j C)\), while the ordinary square \(H^2\)
has lower-left block \(Ca+eC=C\iota_{v+w}\). This proves (HK.9) with the
right-module convention of (HK.4).

Here is the cocycle verification. Associativity gives
(\(D+hH\)) star ((\(D+hH\)) star (\(D+hH\))) equal to
((\(D+hH\)) star (\(D+hH\))) star (\(D+hH\)). The constant and first-order
curvatures are zero, so the coefficient of \(h^2\) is \(DO-OD=0\). The
constant-Poisson product (HK.4) is associative through this order:
\(\exp((h/2)\pi^{ij}\partial_i\otimes\partial_j)\) has commuting differential
operators, and the Leibniz rule makes the two triple products equal.
No global quantized category is being asserted.

There is no endomorphism of cohomological degree one and internal
degree -4: its four block coefficient degrees would be -2,-3,-1,-2.
Consequently any nonzero \(O\) represents a nonzero class in
\(\operatorname{Ext}^2_R(M_b,M_b)\) of internal degree -4.

If \(O\) vanished, (HK.9), evaluated on \(K_2\to K_1\), would give

\[
 \{f_i,C\}_\pi=z_i C\qquad (z_i\in\mathbf C).
 \tag{HK.10}
\]

Because \(\pi\) is nondegenerate and \(C\) is a nonzero linear form, its
Hamiltonian vector field is a nonzero constant vector \(v_C\) with
\(v_C(C)=0\). Choose linear coordinates with \(C=x\) and \(v_C=\partial_y\).
Equation (HK.10) says every homogeneous quadric \(f_i\) has \(y\)-dependence
only in the monomial \(xy\). All four quadrics therefore vanish at the
nonzero projective point on the \(y\)-axis, contradicting regularity.
Hence \(O\) is nonzero. This argument does not require (and does not
assert) the existence of any second-order deformation.

Other internal degrees of a first-order correction cannot affect
the weight -4 equation: degree -3 is zero, all remaining degrees
are at least -2, and the coefficient of the first Moyal correction
lowers internal degree by two. The second-order correction has no
degree -4 part. This exhausts all corrections for the exactly
homogeneous module in (HK.2).

\subsection{4. The ordinary local-to-global injection}\label{SF-E846:s4}

Let \(E\) be a coherent sheaf on a smooth proper fourfold \(X\) with locally
free hull \(H\) and finite kernel and cokernel of \(\phi:E\to H\). Put
\(\mathcal C=\operatorname{cone}(\phi)\), in the distinguished triangle

\[
 E\xrightarrow\phi H\longrightarrow\mathcal C
      \xrightarrow\delta E[1].
 \tag{HK.11}
\]

All cohomology of \(\mathcal C\) is supported on a finite set \(Z\). Since \(E\)
is perfect, \(R\operatorname{Hom}(E,\mathcal C)\) also has finite coherent cohomology
supported on \(Z\). Its support spectral sequence has no rows of
positive sheaf cohomology. Restriction and completion consequently
give canonical isomorphisms

\[
 \operatorname{Ext}^j_X(E,\mathcal C)
   \simeq\bigoplus_{p\in Z}
       \operatorname{Ext}^j_{\widehat{\mathcal O}_{X,p}}
            (M_p,C_p).
 \tag{HK.12}
\]

For any \(M_p\) with a free hull and finite kernel/cokernel, the two
short exact sequences through the hull image, together with
\(\operatorname{Ext}^j_R(\text{finite length},R)=0\) for \(j<4\), give
\(\operatorname{Ext}^1_R(M_p,R^r)=\operatorname{Ext}^2_R(M_p,R^r)=0\). Therefore the local triangle gives
\(\operatorname{Ext}^1_R(M_p,C_p)\simeq\operatorname{Ext}^2_R(M_p,M_p)\).

The actual preceding global arrow \(\operatorname{Ext}^1_X(E,H)\to\operatorname{Ext}^1_X(E,\mathcal C)\)
is zero. This is not inferred from a spectral-sequence edge alone.
The commutative square has top row that arrow, bottom row
the map from the direct sum of \(\operatorname{Ext}^1_R(M_p,R^r)\) to the direct sum
of \(\operatorname{Ext}^1_R(M_p,C_p)\), and vertical arrows restriction followed by
completion. The bottom source is zero, while the right vertical
arrow is the isomorphism (HK.12). Hence the top arrow is zero.
The global triangle then gives
a canonical injection

\[
 j:\bigoplus_{p\in Z}\operatorname{Ext}^2_R(M_p,M_p)
       \hookrightarrow\operatorname{Ext}^2_X(E,E).
 \tag{HK.13}
\]

It is the connecting map after the identifications just made and is
a section of ordinary restriction to the local Ext groups. It is
natural for isomorphisms of the entire hull triangle, in particular
for every specified coherent translation/tensor action.

\subsection{5. The trace-weight lemma, including the canonical volume}\label{SF-E846:s5}

Suppose \(X\) is an abelian fourfold, and at \(p\) the hull map is the graded
map \(M_b\to R\) of section 1, with finite direct-sum multiplicities and
possibly a free summand. Let \(u\) be a homogeneous local \(\operatorname{Ext}^2\) class of
degree \(d\), and let \(a\in\operatorname{Ext}^1_R(M_b,C_p)\) be its preimage. The connecting
map has internal degree zero. For a global degree-two Hochschild
class \(\gamma\), naturality of composition gives

\[
 \operatorname{Tr}_E\bigl(j(u)\,\rho_E(\gamma)\bigr)
   =\operatorname{Tr}_E\bigl(\delta\,[a\,\rho_E(\gamma)]\bigr),
 \tag{HK.14}
\]

up to the single harmless cyclic sign. The bracketed term belongs
to \(\operatorname{Ext}^3_X(E,\mathcal C)\), so (HK.12) makes it depend only on the actual
completed restriction of \(\rho_E(\gamma)\). In particular the \(H^1(T_X)\)
and \(H^2(\mathcal O_X)\) HKR components restrict to zero; their possible global
hypercohomology representatives contribute no extra boundary term
to this equality in the derived category. The remaining constant
bivector evaluation is homogeneous of internal degree -2.

Fix a nonzero invariant volume \(\Omega\) on \(X\). It identifies the
Calabi--Yau trace used on the left of (HK.14) with the Serre trace
after tensoring by \(\Omega\). Rotate this latter trace from \(E\) to the
finite-supported perfect object \(\mathcal C\). This is the usual cyclic
identity for evaluation and coevaluation of bounded locally free
complexes; its sign depends only on the two cohomological degrees.
Thus the result is, up to sign, the Serre trace on \(\mathcal C\) of
the composition of \(a\rho\) with \(\delta\otimes\Omega\).

The finite-support comparison can be made directly at the complex
level. The complex \(R\operatorname{Hom}(C_p,C_p\otimes\omega_R)\) has finite-supported
cohomology, so applying \(R\Gamma_{\mathfrak m}\) to it is an isomorphism. Apply \(R\Gamma_{\mathfrak m}\) to
its perfect-complex evaluation/supertrace map to \(\omega_R\). On degree
four this gives

\[
 \operatorname{Ext}^4_R(C_p,C_p\otimes\omega_R)
   \longrightarrow H^4_{\mathfrak m}(\omega_R)
   \xrightarrow{\operatorname{Res}_p}\mathbf C.
 \tag{HK.15}
\]

For the corresponding global perfect complex the same evaluation
map, followed by forgetting supports, lands in \(H^4(X,\omega_X)\).
These maps commute with restriction and completion. Lipman's
Theorem 9.1 and Corollary 9.2 identify the composite with the global
Serre trace with the last residue in (HK.15): \emph{Dualizing sheaves,
differentials and residues on algebraic varieties}, Astérisque 117
(1984), pp.~82--85.\cite{SFLipman1984} The full invoked proofs, including the comparison
diagrams, were read in this run. The trace/evaluation interpretation
for perfect complexes is also Stacks, Section 48.27 tag 0FVU,
Lemma 48.27.1(5),(7), Remarks 48.27.2--3 and Lemma 48.27.4;\cite{SFStacksE846Duality} its use
here follows by taking the bounded locally free representatives
in the same evaluation diagram, not by extending a scalar formula
without a map of complexes.

For precision, \(\omega_R=R\,dx_1\wedge\cdots\wedge dx_4\). The local-cohomology
group is the span of inverse monomials \(x_1^{-a_1}\cdots x_4^{-a_4}\,dx_1\wedge\cdots\wedge dx_4\),
all \(a_i\ge1\). The residue extracts the coefficient with every \(a_i=1\).
This description is independent of regular formal coordinates;
Lipman §7, pp.~64--69, Lemma 7.3 and Theorem 7.4 give the formula and
its complete coordinate-independence and local-duality proofs.
Consequently the residue is equivariant for simultaneous dilation
of the coordinates and their differentials: it vanishes on every
homogeneous element of nonzero total internal degree. This can
also be checked term by term in the displayed inverse-monomial basis.

The restriction of \(\Omega\) in invariant formal-logarithm coordinates
is a nonzero constant multiple of \(dx_1\wedge\cdots\wedge dx_4\), hence
has degree +4. The connecting map \(\delta\) has degree zero with the
graded hull map of §1. For \(u\) of degree \(d=-4\), the first factor \(a\rho\)
has degree -6, while \(\delta\otimes\Omega\) has degree +4. The value
in (HK.15) therefore has total degree -2, so its residue is zero.
The same argument works for any \(d\) with \(d+2\) nonzero.

Thus the image \(j(u)\) pairs to zero with every ordinary Hochschild
evaluation in degree two. The fourfold Serre trace/evaluation
adjunction supplied and checked in DEG\_L S20/S21 identifies this
annihilator with the kernel of the full semiregularity map \(\sigma^2\).
This invokes precisely the adjunction, as well as the nondegenerate
Hochschild pairing on an abelian fourfold; the first-order identity
\(\sigma\rho=\operatorname{contr}(\operatorname{ch})\) alone would not suffice.

For multiple copies of \(M_b\), take \(u\) tensor the identity on the
multiplicity space and zero on a free summand. Naturality of
Hochschild evaluation makes its restriction block diagonal under
any actual module isomorphism. The trace on these copies is the
sum of the same zero pairings. The free summand's hull cone is
acyclic, so it contributes nothing to (HK.15).

\subsection{6. Actual finite actions, with the stabilizer and scalar cocycle}\label{SF-E846:s6}

Assume the completed local module is actually isomorphic to
\(M_b\otimes\mathbb C^m\) direct-sum \(R^s\), with \(m>0\). This is an isomorphism of
modules, not merely of their associated graded modules or \(K\)-classes.
The hull map is functorial under it, since double duals are functorial.
We prove the existence of an invariant nonzero class without
assuming the original action preserves a grading or chosen resolution.

First consider only the point stabilizer \(G_p\) in the specified finite
translation/tensor action. If a translation \(t_x\) fixes \(p\), then \(x=0\).
Therefore \(G_p\) acts trivially on the completed base ring; the action
at this stage is \(R\)-linear, not semilinear. Its Picard line bundles
are torsion and have flat local trivializations. Coherent global
product isomorphisms between these flat line bundles are constant
in such trivializations: their ratios are global holomorphic
functions on the connected proper abelian variety. Hence the
local operators satisfy \(gh=\alpha(g,h)(gh)\) for one constant scalar
cocycle \(\alpha\) in \(Z^2(G_p,\mathbb C^*)\). Scalar ambiguities disappear on Ext.

The twisted group algebra \(\mathbb C_\alpha[G_p]\) is semisimple. Indeed if a
subspace has a linear projection onto it, averaging its conjugates
over \(G_p\) makes an invariant projection; the scalar factors in
conjugation cancel. Its finite-dimensional simple representations
\(S_\lambda\) therefore give an \(R\)-linear isotypic decomposition

\[
 M_p\simeq\bigoplus_\lambda S_\lambda\otimes N_\lambda,
 \qquad N_\lambda=\operatorname{Hom}_{\mathbf C_\alpha[G_p]}
                    (S_\lambda,M_p).
 \tag{HK.16}
\]

The modules \(N_\lambda\) are finite over \(R\), since they are images of
idempotents on finite direct sums of \(M_p\). The action on the right
is on \(S_\lambda\) alone.

We spell out the required local decomposition fact. Every
endomorphism of \(M_b\) induces multiplication by a scalar \(a\in R\) on
its rank-one double dual and multiplication by \(t\) in \(A\) on its torsion
\(A(-3)\). The extension-class condition is
\((t-a\bmod I)b_{ij}=0\) in \(A\) for all \(i\),\(j\). This follows by applying the
two induced endomorphisms to the actual class in \(\operatorname{Ext}^1_R(I,A)\),
which the Koszul resolution identifies with \(A\otimes\Lambda^2U^*\).
Since one \(b_{ij}\) is a nonzero linear class, this condition forces
\(t(0)=a(0)\). If \(a(0)\) is nonzero, both the induced torsion and hull-image
maps are invertible, hence so is the original endomorphism by the
short exact sequence (HK.2). If \(a(0)=0\), its difference from the
identity is invertible by the same argument. Thus \(\operatorname{End}_R(M_b)\) is
local, with residue \(\mathbb C\), and \(M_b\) is indecomposable. It is not \(R\):
its torsion \(A\) is nonzero. The ring \(\operatorname{End}_R(R)=R\) is local as well.

For completeness, the Krull--Schmidt assertion used here holds
for finite modules over a complete Noetherian local ring. Their
endomorphism algebras are finite and complete over \(R\). Modulo the
maximal ideal they are finite-dimensional algebras; primitive
idempotents lift through their nilpotent radicals and then through
the powers of the maximal ideal. At the complete step, the usual
Newton correction to \(e^2-e\) gives a convergent lift, since \(2e-1\) is
invertible near an idempotent and commutes with every polynomial
in \(e\). This gives a finite decomposition into summands with local
endomorphism rings. Uniqueness follows by writing the identity
of any indecomposable summand as the sum of its composites through
another decomposition. At least one composite is a unit in its
local endomorphism ring. It identifies that summand with one on
the other side; split it off and repeat. In particular every direct
summand of \(M_b^m\) direct-sum \(R^s\) is a sum of copies of \(M_b\) and \(R\).

Apply this fact to all \(N_\lambda\) in (HK.16). There is an actual
\(G_p\)-equivariant isomorphism

\[
 M_p\simeq(M_b\otimes V)\oplus(R\otimes W),
 \tag{HK.17}
\]

where \(V\) and \(W\) carry projective representations with the same
cocycle \(\alpha\), and \(\dim V=m>0\). This isomorphism need not preserve
the original grading; the homogeneous grading is used on its
\(M_b\) factor after the isomorphism, and all characteristic maps
and pairings are transported naturally. Thus \(u\otimes\operatorname{id}_V\)
is a nonzero \(G_p\)-invariant local \(\operatorname{Ext}^2\) class, and §5 proves that
its global image has zero semiregularity.

Finally translate this class over the \(G\)-orbit of \(p\). Choose one
representative for each coset of \(G_p\) and transport by the actual
coherent functor. The result is independent of the representatives,
because the starting class is \(G_p\)-invariant and the action on Ext
is honest. Its direct sum over the orbit is nonzero and \(G\)-invariant.
The injection (HK.13) is natural for the hull triangle, so the
global image remains nonzero and invariant. Each translated
trace pairing is zero by naturality, or by the same formal-log
calculation at that point. Hence it lies in \(\ker\sigma^2\) on invariant Ext.

\subsection{7. The balanced-character consequence and its exact hypotheses}\label{SF-E846:s7}

Let \(X\) be an abelian fourfold of genuine real-multiplication type
(2,2), with polarization \(\theta=\theta_++\theta_-\). Let \(E\) be a coherent
sheaf whose canonical hull map has locally free hull and finite
kernel/cokernel. Suppose at one point its completed module, in
invariant formal-logarithm coordinates, has the exact type
\(M_b^m\) direct-sum \(R^s\) of §6. Suppose moreover

\[
 e^B\operatorname{ch}(E)=r v_q,
 \qquad
 v_q=(1-q\theta_+^2/2)(1-q\theta_-^2/2),
 \qquad r>0,\quad q>0,
 \tag{HK.18}
\]

where \(B\) is any rational algebraic divisor. Then \(E\) is not weakly
\(G\)-semiregular for any specified effective coherent finite
translation/tensor action as in supplied S21.

Suppose to the contrary that it is. The two cap-kernel classes
\(\xi_\pm=\pi_\pm+q\eta_\pm\) of E819§4, conjugated by the full \(B\)-field shear,
retain complementary constant bivector projections \(\pi_\pm\). Their
sum \(\xi\) has nondegenerate projection \(\pi=\pi_++\pi_-\). Evaluation lies
in invariant Ext by supplied S21, and the published
Buchweitz--Flenner composition identity gives \(\sigma^2\rho^2(\xi)=0\).
Injectivity of \(\sigma^2\) on that invariant Ext forces \(\rho^2(\xi)=0\).
Restrict and complete. The \(H^1(T)\) and \(H^2(\mathcal O)\) components vanish,
so the local module has a flat first-order deformation for \(\pi\),
by the explicit bar/free-resolution argument of E819§3.

Under any actual module splitting (HK.17), the characteristic
class restricts diagonally to each \(M_b\) copy. Thus \(M_b\) itself has
zero first-order characteristic class and a first-order lift.
Sections 2--3 construct a nonzero local class \(u\) of degree -4.
Sections 4--6 give a nonzero invariant global \(\operatorname{Ext}^2\) class in \(\ker\sigma^2\),
contradiction. The second-order local obstruction was used only
as an explicit cocycle construction in ordinary Ext. At no point
is a second-order noncommutative lift deduced from semiregularity.

The published first-order input is Buchweitz--Flenner,
\emph{The global decomposition theorem for Hochschild (co-)homology of
singular spaces via the Atiyah--Chern character}, Adv. Math. 217 (2008),
Theorem 6.4.2, with Proposition 6.2.1, Proposition 6.3.1, Corollary 6.3.2
and the proof through Proposition 6.4.4 (author-copy pp.~29--33).\cite{SF-BF2008}
The invariant-evaluation and CY4-adjunction interfaces are the
explicit supplied proved inputs DEG\_L S20/S21: S20's first proof
paragraph (the Serre-transpose identity), S21.6 and the paragraphs
immediately preceding S21.8. Those exact passages were reread in
this run; no publication status is added to the supplied results.

\subsection{8. Actual global objects in the stated local scope}\label{SF-E846:s8}

The local type is an actual coherent-sheaf type, not only a
formal character. For any regular homogeneous four-quadric
sequence the Hilbert series of \(A\) is \((1+t)^4\), so \(\mathfrak m^5\) is contained
in \(I\). The quotient \(R\to A\) therefore factors through \(R/\mathfrak m^5\). In
invariant formal-logarithm coordinates at a point of an abelian
fourfold, the canonical isomorphism of finite jets with the
algebraic local ring realizes it as a closed algebraic subscheme
\(Z\) of length 16. Its completed ideal is exactly \(I\).

There is a global extension
\(0\to\mathcal O_Z\to M\to\mathcal I_Z\to0\) with any prescribed local class \(b\). Indeed
\(R\operatorname{Hom}(\mathcal I_Z,\mathcal O_Z)\) has finite-supported cohomology, and its support
spectral sequence collapses, so global \(\operatorname{Ext}^1\) identifies with the
completed local \(\operatorname{Ext}^1\). The Koszul calculation identifies the latter
with \(A\otimes\Lambda^2U^*\), in which the six specified linear classes
are legitimate coordinates. The resulting sheaf has local module
\(M_b\), canonical hull \(\mathcal O_X\), and identical length 16 kernel and
cokernel. Additivity gives \(\operatorname{ch}(M)=\operatorname{ch}(\mathcal O_X)=1\).

For a concrete balanced ambient character use the actual rank 32
polystable bundle \(H\) of E815/E820 on the CM RM fourfold \(E_i^4\).
Its decomposition into actual line bundles has multiplicities
\(24,2,2,1,1,1,1\) and \(\operatorname{ch}(H)=32v_{1/2}\). Distinct line types have no
Hom between them: their normalized differences are the primitive
index-two classes checked in E813/E815. Thus every global hull
endomorphism is a constant block matrix on the multiplicity
spaces, with algebra \(\operatorname{Mat}_{24}(\mathbb C)\) direct-sum \(\operatorname{Mat}_2(\mathbb C)\) direct-sum
\(\operatorname{Mat}_2(\mathbb C)\) direct-sum \(\mathbb C^4\).

Choose 33 distinct points \(p_1,\ldots,p_{33}\). At each choose the same
quadratic CI in invariant logarithm coordinates and its length 16
algebraic subscheme \(Z_j\). In local frames of the 32 line summands,
define a surjection \(H\to Q=\bigoplus_j\mathcal O_{Z_j}\). At \(p_j\) for \(j\le32\)
use the row \(e_j\); at \(p_{33}\) use the row \((1,\ldots,1)\). Each is a surjection,
and arbitrary such jets define actual morphisms to the finite
sheaf \(Q\). Let \(K\) be its kernel and \(T=Q\). At each point the hull
sequence has \(K_p\simeq I\) direct-sum \(R^{31}\). The support spectral sequence
again identifies \(\operatorname{Ext}^1_X(K,T)\) with the direct sum of its completed
local groups. Choose the nonzero scalar linear class \(b\) on the
\(I\) factor at every point. All \(\operatorname{Ext}^1\) from the free local factors
vanish, so this constructs an actual coherent extension

\[
 0\longrightarrow T\longrightarrow E\longrightarrow K
       \longrightarrow0,
 \qquad E_p\simeq M_b\oplus R^{31}\quad(p\in\{p_j\}).
 \tag{HK.19}
\]

Its maximal torsion is \(T\), its double dual is \(H\), and its hull
cokernel is \(Q\). Hence its full character is \(\operatorname{ch}(H)\). It is an
ordinary gluable object because it is a coherent sheaf.

This \(E\) is indecomposable. Let \(e\) be an idempotent endomorphism.
Its induced hull endomorphism \(a\) preserves \(K\) and hence induces
an endomorphism on \(Q\). At \(p_j\), \(j\le32\), compatibility with the
quotient row \(e_j\) forces the \(j\)-th row of the constant block
matrix \(a\) to be a scalar multiple of \(e_j\). Thus \(a\) is diagonal.
The last row forces all 32 diagonal entries to agree, so
\(a=\lambda\operatorname{id}\). Since \(e\) is idempotent, \(\lambda\) is 0 or 1.
If \(\lambda=0\), the induced endomorphism of each torsion algebra \(A\)
is multiplication by an idempotent \(t_j\). Compatibility with \(b\)
gives \(t_jb=0\), so \(t_j\) has zero residue, as in §6. A local Artin
algebra has no nonzero idempotent of zero residue, hence \(t_j=0\).
Now \(e\) has image in \(T\) and vanishes on \(T\), so \(e^2=0\) and therefore \(e=0\).
If \(\lambda=1\), apply this argument to \(1-e\) to get \(e=1\). There is no
nontrivial idempotent, which proves indecomposability.

The identity action is an actual finite coherent action on these
examples; no extra action is presumed to exist. The theorem
excludes weak semiregularity for it and for every further specified
action if one exists. This is actual-family evidence, not a strict
comparison with all earlier local Artin criteria. No particular
quadratic CI surviving those earlier criteria is asserted here.

\subsection{9. Current boundary and certificate scope}\label{SF-E846:s9}

There is no claim here for arbitrary filtered quadratic jets, arbitrary
matrix-valued \(b\), nonfree Artin defects, or all complete intersections.
No second-order noncommutative semiregularity theorem is imported.
The failed broad bridge through Perry remains separate: his 2026
preprint Theorem 4.1 assumes geometric origin. Pridham's published
2024 cyclic argument has explicit additional hypotheses in section 1.5
and Remark 1.21; these have not been verified for the present Moyal-plus-
gerbe deformation, and are not used in the proposed ordinary trace route.

The entire local calculation, support/trace interfaces and finite-action
step have passed independent whole-proof review. A strict scope comparison
with all previous Artin exclusions has not been established; therefore no
independent main-gate movement is claimed.

The exact author checker rm\_homogeneous\_koszul\_certificate.py
checks the complete finite exterior operator spaces: the lower
chain space has dimension 1 (69 variables, rank 68), the top
anticommutant dimension 4 (56 variables, rank 52), and the truncated
bottom anticommutant dimension 4 (52 variables, rank 48). It also
checks the degree bound, the Hamiltonian-axis implication, and
the entire 31-by-31 second-order matrix identity at three rational
inputs, retaining the first-order Bianchi remainder. The first
actual run passed 34 checks on 2026-11-22 at 02:13:46--02:13:51UTC,
elapsed 5.900344789 seconds. These arithmetic checks neither produce
a first-order existence witness nor verify the global argument.
